% ==========================================================================
% Capítulo: Metodología
% Estado: se registra en ../STATUS.md, no aquí.
%
% Debe contener (project-context/formato-anteproyecto.md, "Metodología"):
%   - Formulación y JUSTIFICACIÓN de la metodología de ingeniería elegida,
%     soportada en referentes de la profesión (marcos, estándares, modelos
%     de proceso) y citada. Explicar el "cómo" y por qué esa y no otra.
%   - Estructura lógica del proceso: recolección de datos, organización,
%     sistematización, análisis, interpretación y presentación de
%     resultados.
%   - FASES DE DESARROLLO DEL PROYECTO: en estrecha relación con los
%     objetivos específicos; cada fase con duración asignada.
%   - ANÁLISIS DE RIESGOS Y LIMITACIONES: riesgos que afectan alcance,
%     cronograma, presupuesto y calidad.
%   - CRONOGRAMA (diagrama de Gantt).
% NO debe contener: metodología descrita sin justificación, ni fases que no
% mapeen a un objetivo específico.
%
% Hechos técnicos: ../../project-context/technologies.md y requirements.md.
% ==========================================================================
\chapter{Metodología}\label{cha:metodologia}

Este capítulo formula y justifica la metodología de ingeniería adoptada para resolver el problema planteado en el \cref{cha:descripcion-problema}, describe la estructura lógica del proceso de investigación, desde la recolección de datos hasta la presentación de resultados, define las fases de desarrollo del proyecto en relación directa con los objetivos específicos del \cref{cha:objetivos}, presenta el cronograma del proyecto y analiza los riesgos y limitaciones que pueden afectar su alcance, cronograma y calidad.

\section{Estrategia metodológica}\label{sec:estrategia-metodologica}

El proyecto adopta un proceso de desarrollo iterativo e incremental de tipo Scrum, siguiendo la \emph{Scrum Guide} de \citet{schwaber-scrumguide-2020}, organizado en \emph{sprints} de tres semanas. Esta elección responde a que el alcance del proyecto, cuatro objetivos específicos con dependencias técnicas parciales entre sí, se beneficia de ciclos cortos de construcción y retroalimentación, en los que cada \emph{sprint} entrega un incremento verificable y permite reordenar el trabajo restante conforme avanza el desarrollo y conforme lo permite la disponibilidad de los nodos de cómputo del IAsLab, en lugar de comprometer de antemano una secuencia rígida de actividades como exigiría un modelo de desarrollo en cascada.

La duración de tres semanas por \emph{sprint} responde a la naturaleza del trabajo de las primeras iteraciones. Levantar el clúster comprende una cadena de actividades, entre ellas desplegar Kubernetes sobre las estaciones de trabajo de la sala 104M, configurar la red, el nodo de control y los nodos de trabajo, aprovisionar el acceso a las unidades de procesamiento gráfico y, sobre esa base ya estable, integrar las extensiones que la plataforma necesita. Cada una de ellas depende de que la anterior quede estable, de modo que una iteración más corta obligaría a cerrar \emph{sprints} sin un clúster verificable. El equipo adoptó por eso una cadencia de tres semanas, que permite terminar cada iteración con infraestructura funcionando y comprobable.

% [verify: ADR-030: Jira y Bitbucket salen del acta 2026-08-26 (Compromisos), no de project-context]
Los artefactos y ceremonias que el equipo emplea de forma efectiva son los que exige el curso de proyecto de grado, es decir, un repositorio Jira para el \emph{backlog} y el seguimiento de las iteraciones, un repositorio Bitbucket para el control de versiones y los archivos de configuración de la infraestructura, y una reunión semanal de seguimiento con el tutor del proyecto, en la que se revisa el avance de la iteración en curso y se ajusta el alcance de la siguiente. El curso de Proyecto de Grado I exige, como mínimo, el anteproyecto y dos \emph{sprints}. No se documentan otras ceremonias de Scrum, por ejemplo el \emph{daily scrum} o las retrospectivas formales, porque el equipo no las practica de forma sistemática; el punto de control periódico del proyecto es la reunión semanal con el tutor.

La estructura lógica del proceso de investigación se organiza en cuatro etapas. La recolección de datos se apoya en tres fuentes instrumentadas por el propio sistema, que son la telemetría de \emph{hardware} e inferencia que produce el módulo de observabilidad, los resultados que genera el motor de \emph{benchmarking} automatizado sobre los nodos de inferencia, y los instrumentos de usabilidad aplicados a los usuarios finales de la plataforma. La organización y sistematización de estos datos consiste en consolidar las métricas de \emph{hardware}, de inferencia y de percepción de usuario en formatos comparables entre fases y entre motores de inferencia evaluados. El análisis contrasta estas métricas contra los objetivos específicos formulados, determinando en qué medida cada módulo construido resuelve la carencia que motivó su desarrollo. Finalmente, la interpretación y presentación de resultados traduce los hallazgos del análisis en conclusiones sobre el nivel de aceptación operacional de la plataforma.

La evaluación empírica prevista en \cref{obj:evaluacion} se apoya en dos instrumentos con respaldo metodológico propio y en un criterio cuantitativo de aceptación, según el cual el sistema debe soportar 20 usuarios simultáneos utilizando la plataforma. Para el desempeño técnico, el proyecto recurre a pruebas de carga sobre los nodos de inferencia bajo esa concurrencia objetivo; \citet{jiang-loadtestingsurvey-2015}, en su revisión de la literatura sobre pruebas de carga en sistemas de gran escala, sistematizan el diseño de la carga, su ejecución y el análisis de los resultados obtenidos, estructura que este proyecto sigue para contrastar el comportamiento de los motores de inferencia desplegados bajo esa concurrencia realista de usuarios. Para la aceptación operacional por parte de los usuarios finales, se emplea el \emph{System Usability Scale}, un cuestionario de diez ítems utilizado como instrumento estándar para cuantificar la usabilidad percibida de un sistema de \emph{software}. \citet{lewis-susbenchmarks-2018} recalibraron sus bandas de referencia por ítem a partir de 446 estudios y más de cinco mil respuestas individuales, lo que permite interpretar el puntaje obtenido frente a una línea base publicada en lugar de valorarlo de forma aislada. El instrumento se aplica a los distintos roles de usuario del IAsLab identificados en el \cref{cha:motivacion-antecedentes}.

\section{Fases de desarrollo del proyecto}
\label{sec:fases-desarrollo}

Las fases de desarrollo del proyecto guardan correspondencia directa con los objetivos específicos formulados en el \cref{cha:objetivos} y se estructuran según una cadena estricta de dependencias técnicas. El proceso inicia con la puesta en marcha del clúster sobre las estaciones de la sala 104M y el despliegue de los motores de inferencia base. Sobre esa infraestructura se instrumenta el módulo de observabilidad, necesario para recolectar telemetría y métricas operativas antes de aplicar restricciones de recursos. Posteriormente, se implementa la gobernanza de cuotas dentro del clúster y se construye la plataforma web de despliegue y administración, la cual expone y gestiona las capacidades previamente estabilizadas. El ciclo culmina con la evaluación empírica del sistema integrado mediante pruebas de carga y valoración de usabilidad. La \cref{tab:fases-desarrollo} sintetiza esta secuencia, detallando las actividades principales de cada fase, su articulación con los objetivos específicos y su duración estimada en \emph{sprints} de tres semanas.

\begin{xltabular}{\textwidth}{p{0.21\textwidth} X p{0.14\textwidth} p{0.09\textwidth}}
        \caption{Fases de desarrollo del proyecto, su correspondencia con los objetivos específicos y su duración.}
        \label{tab:fases-desarrollo} \\
        \toprule
        \textbf{Fase} & \textbf{Actividades principales} & \textbf{Objetivo específico} & \textbf{Duración} \\
        \midrule
        \endfirsthead
        \toprule
        \textbf{Fase} & \textbf{Actividades principales} & \textbf{Objetivo específico} & \textbf{Duración} \\
        \midrule
        \endhead
        \bottomrule
        \endfoot
        Puesta en marcha del clúster &
        Despliegue y configuración de Kubernetes sobre las estaciones de trabajo de la sala 104M, aprovisionamiento de las unidades de procesamiento gráfico mediante el NVIDIA GPU Operator, integración de KubeRay y despliegue de los motores de inferencia como imágenes de contenedor &
        \Cref{obj:orquestacion} & 2 \emph{sprints} \\
        \addlinespace
        Observabilidad de infraestructura &
        Instrumentación de métricas de \emph{hardware} y de carga de trabajo por nodo con Prometheus, el exportador DCGM, Loki y Grafana &
        \Cref{obj:telemetria} & 2 \emph{sprints} \\
        \addlinespace
        Gobernanza de cuotas &
        Implementación del esquema de reserva de cupos y de las restricciones de infraestructura dentro del clúster, con prioridad de los profesores sobre la asignación de la capacidad de la sala y acoplamiento a los atributos de rol de SAAMFI &
        \Cref{obj:gobernanza} & 2 \emph{sprints} \\
        \addlinespace
        Plataforma de despliegue y administración &
        Construcción del servicio y la interfaz web con los que un usuario despliega modelos entrenados de forma concurrente, y del módulo administrativo que gestiona las políticas de cuota y presenta la telemetría recolectada &
        \Cref{obj:orquestacion}, \cref{obj:gobernanza} & 3 \emph{sprints} \\
        \addlinespace
        Evaluación empírica &
        Ejecución de pruebas de carga con 20 usuarios simultáneos sobre los nodos de inferencia mediante el motor de \emph{benchmarking} automatizado, y aplicación del \emph{System Usability Scale} con usuarios finales &
        \Cref{obj:evaluacion} & 1 \emph{sprint} \\
\end{xltabular}

\section{Cronograma}
\label{sec:cronograma}

El proyecto inicia el 10 de agosto de 2026 y finaliza en mayo de 2027. Entre el 10 de agosto y el 20 de septiembre de 2026 el equipo dedica una etapa previa a la formulación del anteproyecto, el levantamiento de requisitos con el laboratorio y la definición de la arquitectura propuesta. Los \emph{sprints} de desarrollo comienzan el 21 de septiembre de 2026 y se suceden en ciclos de tres semanas. El formato de anteproyecto de la facultad exige que las actividades del cronograma se deriven de los objetivos específicos del \cref{cha:objetivos} y sean más detalladas y medibles que ellos. La \cref{tab:cronograma} presenta ese desglose, con las fechas de inicio y fin de cada \emph{sprint} y su entregable. La tabla omite la columna de responsable porque todas las actividades las asume el equipo del proyecto en su conjunto.

% [verify: ADR-033: el sprint 4 termina el 13-dic, y el sprint 5 (11-31 ene) inicia en enero]
La ejecución del proyecto se estructura en dos grandes etapas consecutivas. La primera etapa abarca de agosto a diciembre de 2026, periodo en el que se consolida la formulación del anteproyecto y se ejecutan los cuatro primeros \emph{sprints} de desarrollo, correspondientes a la puesta en marcha del clúster de cómputo y a la instrumentación de observabilidad de infraestructura. La segunda etapa comprende de enero a mayo de 2027, periodo durante el cual se implementan el sistema de gobernanza de cuotas, la plataforma de despliegue y administración, y la evaluación empírica del sistema, culminando con la consolidación del informe final y la entrega de la solución terminada.

\begin{xltabular}{\textwidth}{p{0.13\textwidth} X p{0.16\textwidth} p{0.21\textwidth}}
        \caption{Cronograma del proyecto, por \emph{sprints} de tres semanas, con entregables asociados.}
        \label{tab:cronograma} \\
        \toprule
        \textbf{\emph{Sprint}} & \textbf{Actividad} & \textbf{Fechas} & \textbf{Entregable} \\
        \midrule
        \endfirsthead
        \toprule
        \textbf{\emph{Sprint}} & \textbf{Actividad} & \textbf{Fechas} & \textbf{Entregable} \\
        \midrule
        \endhead
        \bottomrule
        \endfoot
        Etapa previa &
        Formulación del anteproyecto, levantamiento de requisitos con el laboratorio y definición de la arquitectura propuesta &
        10 ago -- 20 sep 2026 &
        Documento de arquitectura y especificación de requisitos \\
        \addlinespace
        \emph{Sprint} 1 &
        Despliegue de Kubernetes sobre las estaciones de trabajo de la sala 104M, con configuración de red, nodo de control y nodos de trabajo &
        21 sep -- 11 oct 2026 &
        Clúster Kubernetes operativo y su configuración versionada \\
        \addlinespace
        \emph{Sprint} 2 &
        Aprovisionamiento de las GPU mediante el NVIDIA GPU Operator, integración de KubeRay y despliegue de los motores de inferencia como imágenes de contenedor &
        12 oct -- 1 nov 2026 &
        Modelo desplegado y consumible sobre el clúster \\
        \addlinespace
        \emph{Sprint} 3 &
        Instrumentación de métricas de \emph{hardware} por nodo (temperatura, VRAM, consumo energético, PCIe) y agregación centralizada de registros &
        2 nov -- 22 nov 2026 &
        Prometheus, exportador DCGM y Loki desplegados en el clúster \\
        \addlinespace
        \emph{Sprint} 4 &
        Instrumentación de métricas de inferencia y construcción de los paneles de visualización &
        23 nov -- 13 dic 2026 &
        Panel de telemetría en Grafana (\cref{obj:telemetria}) \\
        \addlinespace
        Hito &
        Entrega y formalización del documento de anteproyecto &
        Inicios de dic. 2026 &
        Documento de anteproyecto \\
        \addlinespace
        \emph{Sprint} 5 &
        Implementación del esquema de reserva de cupos y de las restricciones de recursos dentro del clúster &
        11 ene -- 31 ene 2027 &
        Mecanismo de cuotas operativo en el clúster \\
        \addlinespace
        \emph{Sprint} 6 &
        Acoplamiento de las cuotas a los atributos de rol de SAAMFI y prioridad de los profesores sobre la capacidad de la sala &
        1 feb -- 21 feb 2027 &
        Sistema de gobernanza por rol (\cref{obj:gobernanza}) \\
        \addlinespace
        \emph{Sprint} 7 &
        Construcción del servicio de la plataforma que expone el despliegue de modelos sobre el clúster &
        22 feb -- 14 mar 2027 &
        Servicio de despliegue de modelos \\
        \addlinespace
        \emph{Sprint} 8 &
        Interfaz web de despliegue concurrente de modelos, con selección de nodo y de motor de inferencia &
        15 mar -- 4 abr 2027 &
        Interfaz de despliegue de modelos (\cref{obj:orquestacion}) \\
        \addlinespace
        \emph{Sprint} 9 &
        Módulo administrativo de gestión de políticas de cuota y presentación de la telemetría dentro de la plataforma &
        5 abr -- 25 abr 2027 &
        Módulo administrativo de la plataforma \\
        \addlinespace
        % [verify: ADR-031: diseño propio o reutilización del harness del laboratorio]
        \emph{Sprint} 10 &
        Ejecución del motor de \emph{benchmarking} automatizado bajo carga de 20 usuarios simultáneos y aplicación del \emph{System Usability Scale} &
        26 abr -- 16 may 2027 &
        Reportes de pruebas de carga y de usabilidad (\cref{obj:evaluacion}) \\
        \addlinespace
        Cierre &
        Ampliación del documento del anteproyecto con el desarrollo realizado y los resultados obtenidos, y redacción del informe final &
        17 may -- 31 may 2027 &
        Informe final del proyecto \\
\end{xltabular}

\section{Análisis de riesgos y limitaciones}

El despliegue de una plataforma sobre infraestructura de cómputo compartida conlleva incertidumbres técnicas y de gestión que pueden afectar el alcance, el cronograma o la calidad de la solución. Los riesgos identificados a continuación se derivan de las restricciones operativas y supuestos del proyecto, se concentran en los niveles técnico y metodológico, y se presentan junto con la dimensión que comprometen y las estrategias previstas para su mitigación.

\begin{description}
  \item[Falta de disponibilidad de los nodos de cómputo del IAsLab] Compromete el cronograma, dado que las fases de desarrollo y, en particular, la evaluación empírica dependen del acceso continuo a los nodos de cómputo del laboratorio, condición que el proyecto asume vigente pero no controla directamente. Como estrategia de mitigación, las actividades de desarrollo que no requieren \emph{hardware} de inferencia (por ejemplo, la definición del esquema lógico de cuotas) se priorizan en los periodos de menor disponibilidad de los nodos.
  \item[Cambios o caída del servicio SAAMFI] SAAMFI es el proveedor de identidad institucional de la Universidad Icesi, que emite un \emph{token} con los roles y permisos del usuario, y cada sistema que lo integra decide cómo usar esos atributos. Este riesgo compromete el alcance y el cronograma, pues el control de acceso y la gobernanza de cuotas por rol dependen de los atributos que provee ese servicio. La mitigación consiste en aislar la lógica de mapeo de roles en un componente propio de la plataforma, de modo que un cambio en la interfaz de SAAMFI exija ajustar solo ese componente, sin rediseñar el esquema de cuotas completo.
  % [verify: ADR-030: rango 90-95 % y fuente de poder salen del acta 2026-09-04; FR-03.4 solo dice 90%+]
  \item[Saturación de VRAM y congelamiento de nodos GPU] Compromete la calidad del servicio de inferencia y, potencialmente, el cronograma de las pruebas de carga. El equipo del laboratorio ha confirmado que un uso sostenido de la GPU entre 90\,\% y 95\,\% congela las máquinas de la sala, y que la causa física es que la fuente de poder de la torre no soporta el consumo eléctrico resultante. \citet{liu-gpufailureprediction-2022} documentan que este tipo de fallos es predecible con telemetría suficiente en producción, por lo que la mitigación prevista es que el módulo de observabilidad, formulado en el \cref{obj:telemetria} (\cref{sec:objetivos-especificos}), genere alertas tempranas y active el diagnóstico automatizado de salud de los nodos ante niveles de saturación críticos (definido en los requisitos del sistema para umbrales de VRAM superiores al 90\,\%), antes de que la saturación derive en un congelamiento no controlado del equipo.
  \item[Uso simultáneo de los equipos para actividades académicas] Los equipos de la sala deben seguir siendo usables para las clases mientras el clúster corre sobre ellos, restricción operativa del laboratorio. Compromete la calidad del servicio de inferencia, en tanto la carga docente puede competir por los mismos recursos de cómputo que utiliza la plataforma. La mitigación consiste en que el esquema de cuotas y gobernanza de recursos, formulado en el \cref{obj:gobernanza} (\cref{sec:objetivos-especificos}), reserve capacidad para las cargas del proyecto sin bloquear el uso académico habitual de los equipos, y en ejecutar fuera del horario de clases toda prueba que pueda degradar el rendimiento de las máquinas.
  % [verify: ADR-030: advertencia del tutor y compromiso GitOps salen del acta 2026-09-09]
  \item[Cambios de requisitos que obliguen a rehacer la infraestructura] El tutor advirtió que, al terminar el proyecto, pueden surgir necesidades nuevas que obliguen a reemplazar buena parte de lo construido. Compromete el alcance y el cronograma de mantenimiento posterior al proyecto. La mitigación es el compromiso adquirido por el equipo con GitOps, documentación y automatización de la infraestructura, de modo que un cambio de requisitos permita redesplegar o adaptar el clúster (por ejemplo, en otra sala del laboratorio) sin reconstruirlo desde cero.
  \item[Límites de ancho de banda de la sala] Compromete la calidad de las pruebas de carga, en tanto la red de la sala del IAsLab tiene una capacidad de 10 Gbit/s. El proyecto no requiere una capacidad mayor, porque no contempla inferencia distribuida entre nodos por red. La mitigación consiste en diseñar los escenarios de prueba de carga a partir de esa capacidad real, en lugar de extrapolar cifras de referentes de la literatura obtenidos en infraestructuras de mayor escala.
\end{description}
